%======================================================================
% Chapter 2: Functions and Graphs
%======================================================================
\chapter{Functions and Graphs}\label{ch:functions}

\begin{motivation}
What \emph{is} a machine learning model? Strip away the jargon and the answer
is short: a model is a \emph{function}. It takes an input (an image, a
sentence, the features of a house) and produces an output (a label, a
translation, a price). Learning from data means searching for a function that
produces the right outputs for the inputs we have seen, and, hopefully, for
new inputs too.

Functions are therefore the central objects of this whole subject. In this
chapter we learn what functions are, how to picture them as graphs, and meet
the handful of function families that appear again and again in machine
learning: linear functions, polynomials, exponentials and logarithms, and the
``activation functions'' used inside neural networks. We also learn two ways
of building new functions from old ones, \emph{composition} and
\emph{inversion}. Composition is exactly how deep neural networks are built.
\end{motivation}

\begin{objectives}
\begin{itemize}
  \item explain what a function is, using the words domain, codomain and range,
    and the notation $f(x)$;
  \item plot and read graphs of functions;
  \item work with linear functions, slopes and intercepts, and interpret the
    slope as a rate of change;
  \item solve quadratic equations and find the minimum of a quadratic function;
  \item describe exponential growth and decay and the special number $e$;
  \item use logarithms and their rules, and explain why machine learning
    prefers log-probabilities;
  \item compose functions, find inverse functions, and predict how parameters
    shift and stretch a graph;
  \item recognise the common activation functions: step, sigmoid, tanh, ReLU
    and softplus.
\end{itemize}
\end{objectives}

%----------------------------------------------------------------------
\section{What is a function?}\label{sec:what-is-function}
%----------------------------------------------------------------------

\begin{picturethis}
A function is a machine. You drop an input into the top, the machine applies a
fixed rule, and an output comes out of the bottom. Drop in the same input
twice, and you get the same output twice. A ``squaring machine'' turns $3$
into $9$ and $-4$ into $16$. A ``cat detector'' turns a photograph into one
of the words \emph{cat} or \emph{not cat}.
\end{picturethis}

\begin{figure}[htbp]
\centering
\begin{tikzpicture}[>={Stealth[length=2.5mm]}]
  \node[draw,thick,rounded corners,fill=MLdef!10,minimum width=3.4cm,
    minimum height=1.3cm,align=center] (f) {function $f$\\[-1pt]{\small (a fixed rule)}};
  \draw[->,thick] (-3.6,0) node[left] {input $x$} -- (f.west);
  \draw[->,thick] (f.east) -- (3.6,0) node[right] {output $f(x)$};
  \node[draw,thick,rounded corners,fill=MLexample!10,minimum width=3.4cm,
    minimum height=1.1cm,align=center] (g) at (0,-1.8) {squaring};
  \draw[->,thick] (-3.6,-1.8) node[left] {$3$} -- (g.west);
  \draw[->,thick] (g.east) -- (3.6,-1.8) node[right] {$9$};
  \node[draw,thick,rounded corners,fill=MLml!10,minimum width=3.4cm,
    minimum height=1.1cm,align=center] (h) at (0,-3.4) {image classifier};
  \draw[->,thick] (-3.6,-3.4) node[left] {photo} -- (h.west);
  \draw[->,thick] (h.east) -- (3.6,-3.4) node[right] {``cat''};
\end{tikzpicture}
\caption{A function is a machine that turns each input into exactly one
  output.}
\label{fig:function-machine}
\end{figure}

\begin{definition}[Function]\label{def:function}
A \term{function} $f$ from a set $A$ to a set $B$ is a rule that assigns to
\emph{each} element $x$ of $A$ \emph{exactly one} element of $B$, called
$f(x)$. The set $A$ is the \term{domain} of $f$ (the allowed inputs), and $B$
is the \term{codomain} (where the outputs live). The set of outputs that
actually occur, $\{f(x) : x \in A\}$, is the \term{range} of $f$.
\end{definition}

\begin{notation}
We write $f\colon A \to B$, read ``$f$ is a function from $A$ to $B$''. The
output for input $x$ is written $f(x)$ and read ``$f$ of $x$''. A rule is
often given by a formula, such as $f(x) = x^2$; equivalently
$x \mapsto x^2$, read ``$x$ maps to $x^2$''. In graphs and applications we
often write $y = f(x)$ and call $x$ the \term{input} (or independent
variable) and $y$ the \term{output} (or dependent variable).
\end{notation}

\begin{pitfall}[$f(x)$ is not $f$ times $x$]
The brackets in $f(x)$ do \emph{not} mean multiplication. $f(3)$ means ``the
output of $f$ when the input is $3$''. Likewise $f(a + 1)$ is \emph{not}
$f(a) + 1$: if $f(x) = x^2$, then $f(a+1) = (a+1)^2 = a^2 + 2a + 1$, which is
quite different from $f(a) + 1 = a^2 + 1$. To evaluate $f$ at any input,
replace \emph{every} $x$ in the formula by that input, in brackets.
\end{pitfall}

Two words in the definition deserve emphasis. \emph{Each} input must get an
output, so the formula $1/x$ does not define a function on all of $\R$ (there
is no output for $x = 0$); it defines a function on the domain
$\{x \in \R : x \neq 0\}$. And \emph{exactly one} output: the rule ``$y$ is a
number whose square is $x$'' is not a function, because the input $4$ would
have two outputs, $2$ and $-2$.

\begin{example}[Domains and ranges]\label{ex:domain-range}
Find a suitable domain and the range of (a) $f(x) = x^2$;
(b) $g(x) = \sqrt{x}$; (c) $h(x) = \frac{1}{x-2}$.

\sol (a) Any real number can be squared, so the domain is $\R$. Squares are
never negative, and every non-negative number is a square, so the range is
$[0, \infty)$.

(b) Square roots of negative numbers are not real, so the domain is
$[0, \infty)$; the range is also $[0, \infty)$.

(c) We cannot divide by zero, so $x = 2$ must be excluded: the domain is
$\{x \in \R : x \neq 2\}$. The output can be any non-zero number (to get
output $y$, take $x = 2 + \tfrac1y$), so the range is
$\{y \in \R : y \neq 0\}$.
\end{example}

\subsection{Ways to describe a function}

A function can be described by a formula ($f(x) = 2x + 1$), by a table of
inputs and outputs, by a graph (\cref{sec:graphs}), in words (``double it and
add one''), or by a computer program. A trained neural network is a function
described by a program together with a long list of numbers (its weights).
Often no simple formula exists at all.

\subsection{Functions of several inputs}

Many functions take more than one input. The area of a rectangle,
$A(\ell, w) = \ell w$, takes two. A house-price model might take three:
\[
  \text{price}(\text{size}, \text{bedrooms}, \text{age})
  = 3000\cdot\text{size} + 10\,000\cdot\text{bedrooms} - 500\cdot\text{age}
  + 20\,000 .
\]
In \cref{ch:linalg1} we will bundle the inputs into a single object, a
\emph{vector}, and in \cref{ch:multivariable} we will learn calculus for such
functions. Until then we mostly study functions of one input.

\begin{mlnote}[Learning is searching for a function]
Supervised machine learning starts from examples of input--output pairs
$(x_1, y_1), (x_2, y_2), \dots, (x_n, y_n)$: photos with their labels, houses
with their prices. We imagine that there is some unknown function that
produced the outputs from the inputs, and we try to find a function $f$ with
$f(x_i) \approx y_i$ for all $i$ (``$\approx$'' means ``approximately
equal''). The set of candidate functions we search through, for example
``all straight lines'' or ``all neural networks of a certain size'', is
called the \term{hypothesis space}. Note that a model \emph{must} be a
function: the same input always gives the same output.
\end{mlnote}

\sectionexercises

\begin{exercise}
Which of the following define a function from $\R$ to $\R$? Explain.
(a) $f(x) = 3x - 1$; (b) ``$y$ is a number with $y^2 = x$'';
(c) $f(x) = \tfrac1x$; (d) $f(x) = \abs{x}$.
\end{exercise}
\begin{solution}
(a) Yes. (b) No: for $x > 0$ there are two such $y$, and for $x < 0$ there is
none. (c) Not on all of $\R$, because $x = 0$ has no output; it is a function
on $\{x \in \R : x\neq0\}$. (d) Yes.
\end{solution}

\begin{exercise}
Let $f(x) = x^2 - 3x$. Compute $f(0)$, $f(2)$, $f(-1)$, $f(a)$ and $f(a+1)$.
\end{exercise}
\begin{solution}
$f(0) = 0$, $f(2) = 4 - 6 = -2$, $f(-1) = 1 + 3 = 4$, $f(a) = a^2 - 3a$, and
$f(a+1) = (a+1)^2 - 3(a+1) = a^2 + 2a + 1 - 3a - 3 = a^2 - a - 2$.
\end{solution}

\begin{exercise}
Find the largest possible domain of (a) $\sqrt{x - 2}$; (b) $\tfrac{1}{x+3}$;
(c) $\sqrt{4 - x^2}$.
\end{exercise}
\begin{solution}
(a) Need $x - 2 \ge 0$: $[2, \infty)$. (b) Need $x \neq -3$.
(c) Need $4 - x^2 \ge 0$, i.e.\ $x^2 \le 4$, i.e.\ $\abs{x} \le 2$: $[-2, 2]$.
\end{solution}

\begin{exercise}
Find the range of $f(x) = x^2 + 1$ and of $g(x) = -\abs{x}$, both with domain
$\R$.
\end{exercise}
\begin{solution}
$x^2 \ge 0$, so $f(x) \ge 1$, and every value $\ge 1$ occurs: the range is
$[1, \infty)$. Similarly $-\abs x \le 0$ and every value $\le 0$ occurs: the
range of $g$ is $(-\infty, 0]$.
\end{solution}

\begin{exercise}
A spam filter reads an email and outputs either ``spam'' or ``not spam''.
What are its domain and codomain? Why would it be a problem if the filter
were not a function?
\end{exercise}
\begin{solution}
The domain is the set of all possible emails; the codomain is the two-element
set $\{\text{spam}, \text{not spam}\}$. If it were not a function, the same
email could receive different verdicts, so its behaviour could not be
predicted, tested or relied upon.
\end{solution}

%----------------------------------------------------------------------
\section{Graphs of functions}\label{sec:graphs}
%----------------------------------------------------------------------

To draw a function we use two perpendicular number lines: a horizontal
\emph{$x$-axis} for inputs and a vertical \emph{$y$-axis} for outputs. Every
pair of numbers $(x, y)$ is a point in this plane: go $x$ units to the right
(left if $x$ is negative) and $y$ units up (down if $y$ is negative).
\Cref{ch:geometry} studies this \emph{Cartesian plane} in detail.

\begin{definition}[Graph]\label{def:graph}
The \term{graph} of a function $f$ is the set of all points $(x, f(x))$ with
$x$ in the domain of $f$.
\end{definition}

The simplest way to sketch a graph is to make a table of values, plot the
points and join them with a smooth curve.

\begin{example}[Plotting a graph]\label{ex:plot-quadratic}
Sketch the graph of $f(x) = x^2 - 2x - 3$ for $-2 \le x \le 4$.

\sol A table of values:
\begin{center}
\begin{tabular}{c|ccccccc}
$x$ & $-2$ & $-1$ & $0$ & $1$ & $2$ & $3$ & $4$\\
\hline
$f(x)$ & $5$ & $0$ & $-3$ & $-4$ & $-3$ & $0$ & $5$
\end{tabular}
\end{center}
For instance $f(-2) = 4 + 4 - 3 = 5$. Plotting the points and joining them
gives the U-shaped curve in \cref{fig:quadratic-plot}.
\end{example}

\begin{figure}[htbp]
\centering
\begin{tikzpicture}
\begin{axis}[mlplot,width=9.5cm,height=7cm,xmin=-3.9,xmax=4.9,ymin=-5.2,ymax=6.5,
  xlabel=$x$,ylabel=$y$,xtick={-3,...,4},ytick={-4,-2,...,6}]
  \addplot[PlotBlue,domain=-2.6:4.6,samples=100]{x^2-2*x-3};
  \addplot[only marks,mark=*,mark size=1.7pt,PlotBlue]
    coordinates {(-2,5) (-1,0) (0,-3) (1,-4) (2,-3) (3,0) (4,5)};
  \addplot[only marks,mark=*,mark size=2.4pt,PlotRed]
    coordinates {(-1,0) (3,0)};
  \node[font=\small,PlotRed,anchor=south east] at (axis cs:-1.05,0.15) {root};
  \node[font=\small,PlotRed,anchor=south west] at (axis cs:3.05,0.15) {root};
  \node[font=\small,anchor=west] at (axis cs:1.2,-4.1) {minimum $(1,-4)$};
  \node[font=\small,anchor=east] at (axis cs:-0.15,-3.0) {$y$-intercept $(0,-3)$};
\end{axis}
\end{tikzpicture}
\caption{The graph of $f(x) = x^2 - 2x - 3$. It crosses the $x$-axis at the
  roots $x = -1$ and $x = 3$, crosses the $y$-axis at $-3$, and has its lowest
  point at $(1, -4)$.}
\label{fig:quadratic-plot}
\end{figure}

\subsection{Reading a graph}

A graph shows at a glance much of what we want to know about a function:
\begin{itemize}
  \item the value $f(a)$ is the height of the graph above (or below) the point
    $a$ on the $x$-axis;
  \item the \term[root of a function]{roots} (or \emph{zeros}) are the inputs
    where $f(x) = 0$, i.e.\ where the graph meets the $x$-axis;
  \item the \term{y-intercept} $f(0)$ is where the graph meets the $y$-axis;
  \item $f$ is \term[increasing function]{increasing} where the graph goes up
    as we move right, and \term[decreasing function]{decreasing} where it goes
    down;
  \item the lowest and highest points show the \emph{minimum} and
    \emph{maximum} values.
\end{itemize}
In \cref{fig:quadratic-plot}, $f$ decreases until $x = 1$, then increases; its
minimum value is $-4$, reached at $x = 1$; and $f(x) < 0$ exactly when
$-1 < x < 3$.

\begin{remark}[The vertical line test]
Because a function gives exactly one output for each input, every vertical
line meets the graph of a function at most once. A curve that some vertical
line crosses twice, such as a circle, is not the graph of a function $y=f(x)$.
\end{remark}

\begin{mlnote}[Reading training curves]
The most-viewed graph in machine learning practice is the \emph{training
curve}: the loss (the model's error) plotted against the number of training
steps (\cref{fig:training-curve}). A healthy curve goes down quickly and then
flattens out. A curve that goes up, oscillates wildly or flattens out
immediately tells an experienced practitioner that something is wrong.
Reading such graphs is a daily skill.
\end{mlnote}

\begin{figure}[htbp]
\centering
\begin{tikzpicture}
\begin{axis}[mlplot,axis lines=left,width=9cm,height=5.5cm,xmin=0,xmax=50,
  ymin=0,ymax=2.6,xlabel={training step (thousands)},ylabel={loss},
  xlabel style={anchor=north,at={(axis description cs:0.5,-0.08)}},
  ylabel style={anchor=south,rotate=90,at={(axis description cs:-0.1,0.5)}}]
  \addplot[PlotBlue,domain=0:50,samples=200]
    {0.3 + 2*exp(-x/6) + 0.05*sin(deg(3*x))*exp(-x/40)};
\end{axis}
\end{tikzpicture}
\caption{A typical training curve: the loss falls quickly at first and then
  levels off. The small wiggles come from the randomness of training.}
\label{fig:training-curve}
\end{figure}

\sectionexercises

\begin{exercise}
Make a table of values of $f(x) = 4 - x^2$ for $x = -3, -2, \dots, 3$ and
sketch the graph. Where are the roots, the $y$-intercept, and the maximum?
\end{exercise}
\begin{solution}
Values: $f(\pm3) = -5$, $f(\pm2) = 0$, $f(\pm1) = 3$, $f(0) = 4$. The graph is
an upside-down U. Roots at $x = -2$ and $x = 2$; $y$-intercept $4$; maximum
value $4$ at $x = 0$.
\end{solution}

\begin{exercise}
Use \cref{fig:quadratic-plot} to answer: for which $x$ is
$x^2 - 2x - 3 > 0$? On which interval is the function increasing?
\end{exercise}
\begin{solution}
The graph is above the $x$-axis for $x < -1$ and for $x > 3$. It is
increasing on $[1, \infty)$ (and decreasing on $(-\infty, 1]$).
\end{solution}

\begin{exercise}
Is the circle consisting of all points with $x^2 + y^2 = 1$ the graph of a
function $y = f(x)$? Can you split it into two pieces that are graphs of
functions?
\end{exercise}
\begin{solution}
No: the vertical line $x = 0$ meets it twice, at $(0, 1)$ and $(0, -1)$. The
upper half is the graph of $f(x) = \sqrt{1 - x^2}$ and the lower half is the
graph of $g(x) = -\sqrt{1 - x^2}$, both with domain $[-1, 1]$.
\end{solution}

%----------------------------------------------------------------------
\section{Linear functions}\label{sec:linear-functions}
%----------------------------------------------------------------------

\begin{definition}[Linear function]\label{def:linear-function}
A function of the form
\[
  f(x) = mx + b ,
\]
where $m$ and $b$ are fixed numbers, is called a \term{linear function}. Its
graph is a straight line. The number $m$ is the \term{slope} and $b$ is the
\term{intercept} (the $y$-intercept, since $f(0) = b$).
\end{definition}

\begin{remark}
Strictly speaking, mathematicians call $mx + b$ an \emph{affine} function and
reserve ``linear'' for the case $b = 0$; we will meet this distinction in
\cref{ch:linalg1}. In everyday usage, and in most of machine learning, both
are called linear.
\end{remark}

\subsection{Slope}

The slope measures how steep the line is. Pick any two points
$(x_1, y_1)$ and $(x_2, y_2)$ on the line. The \emph{run} is the horizontal
change $x_2 - x_1$, the \emph{rise} is the vertical change $y_2 - y_1$, and
\[
  m = \frac{\text{rise}}{\text{run}} = \frac{y_2 - y_1}{x_2 - x_1} .
\]
For a straight line this ratio is the same whichever two points we pick.

\begin{figure}[htbp]
\centering
\begin{tikzpicture}
\begin{axis}[mlplot,width=9cm,height=6.5cm,xmin=-1.2,xmax=6.5,ymin=-0.6,ymax=4.6,
  xlabel=$x$,ylabel=$y$,xtick={0,...,6},ytick={0,...,4}]
  \addplot[PlotBlue,domain=-1.2:6.5]{0.5*x+1};
  \draw[PlotRed,thick] (axis cs:1,1.5) -- (axis cs:5,1.5) -- (axis cs:5,3.5);
  \node[font=\small,PlotRed,anchor=north] at (axis cs:3.6,1.45) {run $=4$};
  \node[font=\small,PlotRed,anchor=west] at (axis cs:5.05,2.5) {rise $=2$};
  \addplot[only marks,mark=*,mark size=2pt,black] coordinates {(1,1.5) (5,3.5) (0,1)};
  \node[font=\small,anchor=north west] at (axis cs:0.95,1.45) {$(1, 1.5)$};
  \node[font=\small,anchor=south east] at (axis cs:4.9,3.55) {$(5, 3.5)$};
  \node[font=\small,anchor=south east] at (axis cs:-0.05,1.05) {$b = 1$};
  \node[font=\small,PlotBlue,anchor=south] at (axis cs:2.6,3.4) {$y = \tfrac12 x + 1$};
\end{axis}
\end{tikzpicture}
\caption{The slope is rise over run: here $m = \tfrac{2}{4} = \tfrac12$. The
  line crosses the $y$-axis at the intercept $b = 1$.}
\label{fig:slope}
\end{figure}

\begin{intuition}[Slope is a rate of change]
The slope tells you \emph{how much the output changes when the input increases
by one unit}. If $y = \tfrac12x + 1$, then every step of $1$ to the right
raises $y$ by $\tfrac12$. A positive slope means the line goes up from left to
right, a negative slope means it goes down, a slope of $0$ gives a horizontal
line, and a larger $\abs{m}$ means a steeper line. Slopes carry units: if $x$
is a house size in square metres and $y$ a price in euros, then $m$ is in euros
per square metre. In \cref{ch:calculus}, the idea of slope will be extended
from straight lines to curves; the result is called the \emph{derivative}.
\end{intuition}

\begin{example}[The line through two points]\label{ex:line-through-points}
Find the linear function whose graph passes through $(1, 2)$ and $(4, 11)$.

\sol The slope is $m = \frac{11 - 2}{4 - 1} = \frac93 = 3$. So $f(x) = 3x + b$,
and since $f(1) = 2$ we need $3 + b = 2$, so $b = -1$. The function is
$f(x) = 3x - 1$. \emph{Check:} $f(4) = 12 - 1 = 11$. \checkmark
\end{example}

A convenient alternative is the \emph{point--slope form}: the line with slope
$m$ through the point $(x_1, y_1)$ is
\[
  y - y_1 = m(x - x_1).
\]

\begin{example}[Celsius and Fahrenheit]
The formula $F = 1.8C + 32$ is a linear function of $C$. Its slope $1.8$ says
that each degree Celsius is worth $1.8$ degrees Fahrenheit; its intercept $32$
is the Fahrenheit temperature at $0\,^\circ$C, the freezing point of water.
\end{example}

\subsection{Special lines}

A horizontal line $y = c$ has slope $0$: the output does not depend on the
input at all (a \emph{constant function}). A vertical line $x = c$ is not the
graph of a function (it fails the vertical line test), and its slope is
undefined because the run is $0$. Two lines are \emph{parallel} exactly when
they have the same slope. When $b = 0$, the output is proportional to the
input: doubling $x$ doubles $y$.

\begin{mlnote}[Linear regression]
The simplest machine learning model is \term{linear regression} with one input:
\[
  \hat y = wx + b .
\]
In machine learning the slope is called the \emph{weight} $w$ and the
intercept the \emph{bias} $b$. The weight measures how sensitive the
prediction is to the input: $w = 3000$ euros per square metre means that each
extra square metre adds $3000$ euros to the predicted price. Training finds
the line that best fits a cloud of data points (\cref{ch:linalg2,ch:optimization}).
\end{mlnote}

\sectionexercises

\begin{exercise}
Find the slope and intercept of the line through $(1, 2)$ and $(3, -4)$.
\end{exercise}
\begin{solution}
$m = \tfrac{-4 - 2}{3 - 1} = -3$. Then $2 = -3 \cdot 1 + b$ gives $b = 5$; the
line is $y = -3x + 5$.
\end{solution}

\begin{exercise}
Write the equation of the line with slope $-2$ through the point $(3, 1)$ in
the form $y = mx + b$.
\end{exercise}
\begin{solution}
$y - 1 = -2(x - 3)$, so $y = -2x + 7$.
\end{solution}

\begin{exercise}
Using $F = 1.8C + 32$: (a) by how much does $F$ change when $C$ increases by
$10$? (b) At what temperature do the Celsius and Fahrenheit scales show the
same number?
\end{exercise}
\begin{solution}
(a) By $1.8 \times 10 = 18$ degrees Fahrenheit. (b) Solve $C = 1.8C + 32$:
$-0.8C = 32$, so $C = -40$. Indeed $-40\,^\circ$C $= -40\,^\circ$F.
\end{solution}

\begin{exercise}
A model predicts $\text{price} = 2500 \times \text{size} + 40\,000$. Interpret
the numbers $2500$ and $40\,000$. By how much does the predicted price rise if
the size increases by $12$ square metres? Is the intercept meaningful?
\end{exercise}
\begin{solution}
Each extra square metre adds $2500$ to the predicted price (the slope); the
intercept $40\,000$ is the prediction for a house of size $0$. A size increase
of $12$ raises the prediction by $12 \times 2500 = 30\,000$. The intercept is
not meaningful on its own (there are no houses of size $0$); it just positions
the line so that it fits the data in the realistic range.
\end{solution}

\begin{exercise}
Are the lines $y = 2x + 1$ and $y = 2x - 5$ parallel? Where do the lines
$y = 2x + 1$ and $y = -x + 4$ meet?
\end{exercise}
\begin{solution}
Yes, both have slope $2$ (and different intercepts, so they never meet). For
the second pair, $2x + 1 = -x + 4$ gives $x = 1$ and $y = 3$: they meet at
$(1, 3)$.
\end{solution}

%----------------------------------------------------------------------
\section{Polynomials and quadratic functions}\label{sec:polynomials}
%----------------------------------------------------------------------

\begin{definition}[Polynomial]\label{def:polynomial}
A \term{polynomial} is a function of the form
\[
  p(x) = a_n x^n + a_{n-1}x^{n-1} + \dots + a_1 x + a_0 ,
\]
where $n$ is a whole number $\ge 0$ and the \emph{coefficients}
$a_0, \dots, a_n$ are real numbers with $a_n \neq 0$. The highest power $n$ is
the \term[degree of a polynomial]{degree}.
\end{definition}

Constant functions have degree $0$, linear functions have degree $1$,
\emph{quadratic} functions have degree $2$ (for example $x^2 - 2x - 3$), and
\emph{cubic} functions have degree $3$ (for example $x^3 - 3x$).

\subsection{Quadratic functions}

A \term{quadratic function} has the form $f(x) = ax^2 + bx + c$ with
$a \neq 0$. Its graph is a \term{parabola}: a U shape that opens upwards if
$a > 0$ and downwards if $a < 0$. An upward parabola has a lowest point, and a
downward one a highest point; this turning point is called the \term{vertex}.
The larger $\abs{a}$, the narrower the U.

\begin{figure}[htbp]
\centering
\begin{tikzpicture}
\begin{axis}[mlplot,width=10cm,height=7cm,xmin=-4.5,xmax=5.5,ymin=-4.5,ymax=6.5,
  xlabel=$x$,ylabel=$y$,xtick={-4,...,5},ytick={-4,-2,...,6},
  legend style={font=\small,at={(0.5,-0.08)},anchor=north,legend columns=3,
    draw=none,/tikz/every even column/.append style={column sep=4mm}}]
  \addplot[PlotBlue,domain=-2.55:2.55,samples=100]{x^2};
  \addlegendentry{$y = x^2$}
  \addplot[PlotRed,domain=-3.6:3.6,samples=100]{-0.5*x^2+2};
  \addlegendentry{$y = -\tfrac12x^2 + 2$}
  \addplot[PlotGreen,domain=0.45:5.5,samples=100]{(x-3)^2+1};
  \addlegendentry{$y = (x-3)^2 + 1$}
  \addplot[only marks,mark=*,mark size=2pt,PlotGreen,forget plot] coordinates {(3,1)};
  \node[font=\small,PlotGreen,anchor=north,fill=white,inner sep=1.5pt]
    at (axis cs:3.9,0.7) {vertex $(3,1)$};
\end{axis}
\end{tikzpicture}
\caption{Three parabolas. $y = x^2$ opens upwards with vertex at the origin;
  $y = -\tfrac12 x^2 + 2$ opens downwards and is wider; $y = (x-3)^2 + 1$ is
  the graph of $x^2$ moved $3$ units right and $1$ unit up.}
\label{fig:parabolas}
\end{figure}

\subsection{Solving quadratic equations}

The roots of a quadratic function are the solutions of $ax^2 + bx + c = 0$.
Sometimes we can find them by factoring (\cref{sec:algebra}): since
$x^2 - 2x - 3 = (x - 3)(x + 1)$, the equation $x^2 - 2x - 3 = 0$ becomes
$(x - 3)(x + 1) = 0$. A product is zero only if one of its factors is zero, so
$x = 3$ or $x = -1$, matching \cref{fig:quadratic-plot}.

When factoring is not obvious, there is a formula that always works.

\begin{theorem}[The quadratic formula]\label{thm:quadratic-formula}
The solutions of $ax^2 + bx + c = 0$ (with $a \neq 0$) are
\[
  \keyeq{x = \frac{-b \pm \sqrt{b^2 - 4ac}}{2a}}
\]
The symbol $\pm$ (``plus or minus'') means that there is one solution with $+$
and one with $-$. The number $D = b^2 - 4ac$ is the
\term{discriminant}: if $D > 0$ there are two real solutions, if $D = 0$
there is exactly one, and if $D < 0$ there are no real solutions (the
parabola does not reach the $x$-axis).
\end{theorem}

\begin{example}[Using the quadratic formula]
Solve $2x^2 - 4x - 1 = 0$.

\sol Here $a = 2$, $b = -4$, $c = -1$, so $D = (-4)^2 - 4\cdot2\cdot(-1) = 16 +
8 = 24 > 0$: two solutions,
\[
  x = \frac{4 \pm \sqrt{24}}{4} = 1 \pm \frac{\sqrt{6}}{2},
  \qquad\text{i.e. } x \approx 2.225 \text{ or } x \approx -0.225 .
\]
\end{example}

\subsection{Completing the square: where is the minimum?}

For machine learning, the most important fact about a quadratic is the
location of its minimum. A technique called \term{completing the square}
rewrites $ax^2 + bx + c$ in \emph{vertex form} $a(x - h)^2 + k$, from which the
vertex $(h, k)$ can be read off directly.

\begin{example}[Finding the minimum of a quadratic]\label{ex:complete-square}
Find the minimum value of $f(x) = x^2 - 6x + 10$.

\sol The identity $(x - 3)^2 = x^2 - 6x + 9$ contains the first two terms of
$f$. So
\[
  f(x) = (x^2 - 6x + 9) + 1 = (x-3)^2 + 1 .
\]
A square is never negative, so $(x - 3)^2 \ge 0$, with equality only at
$x = 3$. Hence $f(x) \ge 1$ for every $x$, and the minimum value $1$ is
reached at $x = 3$. The vertex is $(3, 1)$ (the green parabola in
\cref{fig:parabolas}).
\end{example}

In general, the halfway point between the two terms $x^2$ and $bx$ gives the
recipe: $x^2 + bx = \left(x + \tfrac b2\right)^2 - \tfrac{b^2}{4}$. Applying
this to $ax^2 + bx + c$ shows that

\begin{proposition}[Vertex of a parabola]\label{prop:vertex}
The vertex of $f(x) = ax^2 + bx + c$ is at $x = -\frac{b}{2a}$. If $a > 0$,
this is where $f$ attains its minimum; if $a < 0$, its maximum.
\end{proposition}

\begin{remark}
Completing the square on $ax^2 + bx + c = 0$ in general is exactly how the
quadratic formula is derived. You can try it: divide by $a$, complete the
square, and take square roots of both sides.
\end{remark}

\subsection{Higher-degree polynomials}

Polynomials of higher degree can wiggle more: a polynomial of degree $n$ has
at most $n$ roots and at most $n - 1$ turning points
(\cref{fig:cubic}). For large $\abs{x}$, the highest-power term dominates, so
the graph eventually shoots off to $+\infty$ or $-\infty$.

\begin{figure}[htbp]
\centering
\begin{tikzpicture}
\begin{axis}[mlplot,width=9cm,height=6.5cm,xmin=-2.6,xmax=2.6,ymin=-4.2,ymax=4.2,
  xlabel=$x$,ylabel=$y$,xtick={-2,...,2},ytick={-4,-2,...,4}]
  \addplot[PlotBlue,domain=-2.3:2.3,samples=150]{x^3-3*x};
  \addplot[only marks,mark=*,mark size=2pt,PlotRed] coordinates {(-1.732,0) (0,0) (1.732,0)};
  \addplot[only marks,mark=*,mark size=2pt,PlotGreen] coordinates {(-1,2) (1,-2)};
  \node[font=\small,PlotGreen,anchor=south] at (axis cs:-1,2.1) {turning point};
  \node[font=\small,PlotGreen,anchor=north] at (axis cs:1,-2.1) {turning point};
\end{axis}
\end{tikzpicture}
\caption{The cubic $y = x^3 - 3x$ has three roots ($0$ and $\pm\sqrt3$, in red)
  and two turning points (in green).}
\label{fig:cubic}
\end{figure}

\begin{mlnote}[Squared error is a quadratic]
Suppose we want to summarise the three numbers $y_1 = 2$, $y_2 = 4$, $y_3 = 9$
by a single number $c$, and we measure the quality of $c$ by the total squared
error
\[
  L(c) = (2 - c)^2 + (4 - c)^2 + (9 - c)^2 .
\]
Expanding, $L(c) = 3c^2 - 30c + 101$: a quadratic in $c$ with $a = 3 > 0$. By
\cref{prop:vertex} its minimum is at $c = \frac{30}{2\cdot3} = 5$, which is
exactly the mean $\frac{2 + 4 + 9}{3}$. This is no coincidence: the mean is
always the number that minimises the squared error, a fact we will prove with
calculus in \cref{ch:calculus}. More generally, the loss of a linear regression
model is a quadratic function of its parameters, which is why it is so easy to
minimise.
\end{mlnote}

\sectionexercises

\begin{exercise}
Give the degree and leading coefficient (the coefficient of the highest
power) of (a) $4x^3 - x + 7$; (b) $(x + 1)(x - 2)$; (c) $5$.
\end{exercise}
\begin{solution}
(a) Degree $3$, leading coefficient $4$. (b) Expanding gives $x^2 - x - 2$:
degree $2$, leading coefficient $1$. (c) Degree $0$, leading coefficient $5$.
\end{solution}

\begin{exercise}
Solve by factoring: (a) $x^2 - 5x + 6 = 0$; (b) $x^2 - 16 = 0$;
(c) $x^2 + 3x = 0$.
\end{exercise}
\begin{solution}
(a) $(x-2)(x-3) = 0$: $x = 2$ or $3$. (b) $(x-4)(x+4) = 0$: $x = \pm4$.
(c) $x(x+3) = 0$: $x = 0$ or $-3$.
\end{solution}

\begin{exercise}
Use the discriminant to decide how many real solutions each equation has, and
find them: (a) $x^2 + x + 1 = 0$; (b) $x^2 - 6x + 9 = 0$;
(c) $2x^2 - 3x - 2 = 0$; (d) $x^2 + 4x + 1 = 0$.
\end{exercise}
\begin{solution}
(a) $D = 1 - 4 = -3 < 0$: none. (b) $D = 36 - 36 = 0$: one, $x = 3$.
(c) $D = 9 + 16 = 25$: two, $x = \tfrac{3 \pm 5}{4}$, i.e.\ $x = 2$ or
$x = -\tfrac12$. (d) $D = 16 - 4 = 12$: two,
$x = \tfrac{-4 \pm\sqrt{12}}{2} = -2 \pm \sqrt3$, i.e.\ about $-0.268$ and
$-3.732$.
\end{solution}

\begin{exercise}
Complete the square to find the vertex and the minimum or maximum value of
(a) $f(x) = x^2 + 4x + 7$; (b) $g(x) = -x^2 + 6x - 5$.
\end{exercise}
\begin{solution}
(a) $x^2 + 4x + 7 = (x+2)^2 - 4 + 7 = (x+2)^2 + 3$: vertex $(-2, 3)$, minimum
value $3$. (b) $-x^2 + 6x - 5 = -(x^2 - 6x) - 5 = -[(x-3)^2 - 9] - 5 =
-(x-3)^2 + 4$: vertex $(3, 4)$, maximum value $4$.
\end{solution}

\begin{exercise}
For the data $y = (1, 3, 8)$, write the total squared error
$L(c) = \sum_{i=1}^3 (y_i - c)^2$ as a quadratic in $c$, find the value of $c$
that minimises it, and compare with the mean.
\end{exercise}
\begin{solution}
$L(c) = (1-c)^2 + (3-c)^2 + (8-c)^2 = 3c^2 - 24c + 74$. The vertex is at
$c = \tfrac{24}{6} = 4$, which is the mean $\tfrac{1 + 3 + 8}{3} = 4$. The
minimum value is $L(4) = 9 + 1 + 16 = 26$.
\end{solution}

%----------------------------------------------------------------------
\section{Exponential functions}\label{sec:exponential}
%----------------------------------------------------------------------

\begin{motivation}
Some quantities grow by \emph{adding} a fixed amount each step (a salary that
rises by $1000$ per year). Others grow by \emph{multiplying} by a fixed factor
each step (bacteria that double every hour, money earning interest, a rumour
spreading on social media). Multiplicative growth and decay are described by
exponential functions. In machine learning, exponentials turn arbitrary scores
into positive numbers (and so into probabilities), describe how learning rates
are decayed, and explain why signals can explode or vanish in deep networks.
\end{motivation}

\begin{definition}[Exponential function]\label{def:exponential}
For a fixed number $a > 0$ with $a \neq 1$, the function
\[
  f(x) = a^x
\]
is called an \term{exponential function} with base $a$. Thanks to the rules of
\cref{sec:exponents}, $a^x$ makes sense for every real number $x$: for
example $2^{-1} = \tfrac12$ and $2^{1/2} = \sqrt2$.
\end{definition}

\begin{figure}[htbp]
\centering
\begin{tikzpicture}
\begin{axis}[mlplot,width=9.5cm,height=7cm,xmin=-3.2,xmax=3.2,ymin=-0.6,ymax=8.5,
  xlabel=$x$,ylabel=$y$,xtick={-3,...,3},ytick={0,2,...,8}]
  \addplot[PlotBlue,domain=-3.2:3.1,samples=100]{2^x};
  \addplot[PlotRed,domain=-3.1:3.2,samples=100]{0.5^x};
  \addplot[PlotGreen,domain=-3.2:2.13,samples=100]{exp(x)};
  \addplot[only marks,mark=*,mark size=2pt,black] coordinates {(0,1)};
  \node[font=\small,PlotBlue,anchor=west] at (axis cs:2.35,4.3) {$2^x$};
  \node[font=\small,PlotRed,anchor=east] at (axis cs:-2.35,4.3) {$\left(\tfrac12\right)^x$};
  \node[font=\small,PlotGreen,anchor=east] at (axis cs:1.95,7.7) {$e^x$};
  \node[font=\small,anchor=north west] at (axis cs:0.05,0.95) {$(0,1)$};
\end{axis}
\end{tikzpicture}
\caption{Exponential functions. Every graph passes through $(0, 1)$ and stays
  above the $x$-axis. $2^x$ and $e^x$ grow, $(\tfrac12)^x$ decays.}
\label{fig:exponentials}
\end{figure}

The graphs in \cref{fig:exponentials} show the key properties:
\begin{itemize}
  \item $a^x > 0$ for every $x$: an exponential is never zero or negative;
  \item $a^0 = 1$, so every graph passes through $(0, 1)$;
  \item if $a > 1$ the function increases (\emph{exponential growth}); if
    $0 < a < 1$ it decreases (\emph{exponential decay});
  \item $a^{x + 1} = a \cdot a^x$: increasing the input by $1$
    \emph{multiplies} the output by $a$. (For a linear function, increasing the
    input by $1$ \emph{adds} the slope to the output.)
\end{itemize}

\begin{intuition}[Exponential growth eventually beats everything]
Compare $x^2$ with $2^x$:
\begin{center}
\begin{tabular}{c|cccccccccc}
$x$ & $1$ & $2$ & $3$ & $4$ & $5$ & $6$ & $7$ & $8$ & $9$ & $10$\\
\hline
$x^2$ & $1$ & $4$ & $9$ & $16$ & $25$ & $36$ & $49$ & $64$ & $81$ & $100$\\
$2^x$ & $2$ & $4$ & $8$ & $16$ & $32$ & $64$ & $128$ & $256$ & $512$ & $1024$
\end{tabular}
\end{center}
At first they are similar, but after $x = 4$ the exponential races ahead, and
the gap keeps widening: $2^{30}$ is over a billion while $30^2 = 900$. Any
exponential with base greater than $1$ eventually overtakes any polynomial.
\end{intuition}

\subsection{Growth and decay models}

A quantity that starts at $N_0$ and is multiplied by $a$ in every unit of time
is described by $N(t) = N_0 a^t$.

\begin{example}[Doubling and halving]
(a) A colony of $100$ bacteria doubles every hour. How many are there after
$5$ hours? (b) A substance has a \emph{half-life} of $10$ years: every $10$
years, half of it decays. What fraction remains after $30$ years?

\sol (a) $N(t) = 100 \cdot 2^t$, so $N(5) = 100 \cdot 32 = 3200$.
(b) $30$ years is three half-lives, so the remaining fraction is
$\left(\tfrac12\right)^3 = \tfrac18$. As a formula, the amount after $t$ years
is $A_0 \left(\tfrac12\right)^{t/10}$.
\end{example}

\subsection{The number \texorpdfstring{$e$}{e}}

One base is so useful that it has its own name. Suppose a bank pays $100\%$
interest per year on $1$ euro. Paid once at the end of the year, you get
$2$ euros. Paid in two halves of $50\%$, you get $(1 + \tfrac12)^2 = 2.25$.
Paid in $n$ instalments of $\tfrac{100}{n}\%$, you get
$\left(1 + \tfrac1n\right)^n$:
\begin{center}
\begin{tabular}{c|ccccccc}
$n$ & $1$ & $2$ & $4$ & $12$ & $365$ & $10\,000$ & $1\,000\,000$\\
\hline
$\left(1 + \frac1n\right)^n$ & $2$ & $2.25$ & $2.4414$ & $2.6130$ & $2.7146$ &
  $2.71815$ & $2.718280$
\end{tabular}
\end{center}
The more often interest is paid, the more you get, but the amounts approach a
limit: the number
\[
  e = 2.718281828\ldots,
\]
called \term[e (Euler's number)]{Euler's number}. Like $\pi$, it is
irrational. The function $e^x$ is called \emph{the} exponential function,
also written $\exp(x)$. Its special property, which we will see in
\cref{ch:calculus}, is that at every point its slope equals its value. This
makes $e$ the most natural base for calculus, and it is the base used
throughout machine learning.

\begin{mlnote}[Exponentials in machine learning]
Exponentials appear in machine learning in at least three ways.
\begin{itemize}
  \item \textbf{Making things positive.} A neural network produces arbitrary
    real-valued scores, but probabilities must be positive. Since $e^x > 0$ for
    every $x$, applying $e^x$ to the scores makes them positive; this is the
    first step of the \emph{softmax} function (\cref{ch:dl}).
  \item \textbf{Learning-rate decay.} A common schedule multiplies the
    learning rate by a fixed factor after every epoch:
    $\eta_t = \eta_0 \gamma^t$ with, say, $\eta_0 = 0.1$ and $\gamma = 0.95$.
    After $20$ epochs the rate is $0.1 \times 0.95^{20} \approx 0.036$.
  \item \textbf{Exploding and vanishing.} Multiplying many numbers that are
    slightly larger or smaller than $1$ gives exponential behaviour:
    $1.1^{100} \approx 13\,781$ but $0.9^{100} \approx 0.000\,027$. In a deep
    network, signals pass through many layers and get multiplied many times,
    which is why they can \emph{explode} or \emph{vanish}
    (\cref{ch:dl}).
\end{itemize}
\end{mlnote}

\sectionexercises

\begin{exercise}
Evaluate $2^5$, $2^{-3}$, $10^0$, $\left(\tfrac13\right)^{-2}$, $8^{1/3}$ and
$e^0$.
\end{exercise}
\begin{solution}
$32$, $\tfrac18$, $1$, $3^2 = 9$, $2$, $1$.
\end{solution}

\begin{exercise}
A colony of $500$ cells triples every hour. Write a formula for the number of
cells after $t$ hours and compute it for $t = 4$.
\end{exercise}
\begin{solution}
$N(t) = 500 \cdot 3^t$, and $N(4) = 500 \cdot 81 = 40\,500$.
\end{solution}

\begin{exercise}
A substance with a half-life of $8$ days starts at $160$ grams. Write a
formula for the amount after $t$ days, and find the amount after $24$ days.
\end{exercise}
\begin{solution}
$A(t) = 160 \cdot \left(\tfrac12\right)^{t/8}$. After $24$ days (three
half-lives), $A = 160 \cdot \tfrac18 = 20$ grams.
\end{solution}

\begin{exercise}
Compare $x^3$ and $2^x$ at $x = 10$ and $x = 20$.
\end{exercise}
\begin{solution}
At $x = 10$: $x^3 = 1000$ and $2^x = 1024$, nearly equal. At $x = 20$:
$x^3 = 8000$ but $2^x = 1\,048\,576$, more than a hundred times larger.
\end{solution}

\begin{exercise}
A signal is multiplied by $0.5$ in each of $50$ layers of a network. What
fraction of it is left? What if it is multiplied by $1.5$ in each layer?
\end{exercise}
\begin{solution}
$0.5^{50} \approx 8.9\times10^{-16}$: the signal has vanished. And
$1.5^{50} \approx 6.4 \times 10^{8}$: the signal has exploded. Factors only
slightly different from $1$ have enormous effects over many layers.
\end{solution}

%----------------------------------------------------------------------
\section{Logarithms}\label{sec:logarithms}
%----------------------------------------------------------------------

An exponential answers the question ``what do I get if I multiply $a$ by
itself $x$ times?''. A logarithm answers the reverse question: ``how many
times must I multiply $a$ by itself to get $y$?''.

\begin{definition}[Logarithm]\label{def:logarithm}
Let $a > 0$ with $a \neq 1$, and let $y > 0$. The \term{logarithm} of $y$ to
base $a$, written $\log_a y$, is the number $x$ such that $a^x = y$:
\[
  \log_a y = x \quad\text{means exactly}\quad a^x = y .
\]
\end{definition}

So $\log_2 8 = 3$ because $2^3 = 8$; $\log_{10} 1000 = 3$ because
$10^3 = 1000$; $\log_{10} 0.01 = -2$ because $10^{-2} = 0.01$; and
$\log_a 1 = 0$ for every base, because $a^0 = 1$.

\begin{notation}
The logarithm with base $e$ is the \term{natural logarithm}, written $\ln$:
$\ln y = \log_e y$. Base $10$ is common in science and base $2$ in computer
science and information theory. \textbf{In machine learning papers and code,
``$\log$'' with no base almost always means the natural logarithm $\ln$}; for
example, \texttt{np.log} in Python computes $\ln$.
\end{notation}

\subsection{The rules of logarithms}

Because logarithms are exponents, the laws of exponents
(\cref{thm:exponent-laws}) turn into laws of logarithms.

\begin{theorem}[Laws of logarithms]\label{thm:log-laws}
For $x, y > 0$, any real number $p$, and any valid base:
\begin{align*}
  \log(xy) &= \log x + \log y
    && \text{(products become sums)},\\
  \log\left(\tfrac{x}{y}\right) &= \log x - \log y
    && \text{(quotients become differences)},\\
  \log\left(x^p\right) &= p\log x
    && \text{(powers become multiples)},\\
  \log_a x &= \frac{\ln x}{\ln a}
    && \text{(change of base)}.
\end{align*}
\end{theorem}

\begin{proof}[Why the first rule holds]
Let $u = \log_a x$ and $v = \log_a y$, so $x = a^u$ and $y = a^v$. Then
$xy = a^u a^v = a^{u+v}$, which says exactly that $\log_a(xy) = u + v$. The
other rules are proved in the same way from the other exponent laws.
\end{proof}

The first rule is the reason logarithms were invented (in the 1600s): they
turn hard multiplications into easy additions.

\begin{figure}[htbp]
\centering
\begin{tikzpicture}
\begin{axis}[mlplot,width=9.5cm,height=6.5cm,xmin=-0.4,xmax=8.5,ymin=-3.3,ymax=2.6,
  xlabel=$x$,ylabel=$y$,xtick={1,...,8},ytick={-3,...,2}]
  \addplot[PlotBlue,domain=0.04:8.4,samples=200]{ln(x)};
  \addplot[only marks,mark=*,mark size=2pt,black] coordinates {(1,0) (2.71828,1)};
  \node[font=\small,anchor=south east] at (axis cs:0.95,0.08) {$(1,0)$};
  \node[font=\small,anchor=south east] at (axis cs:2.7,1.05) {$(e,1)$};
  \node[font=\small,PlotBlue,anchor=south] at (axis cs:6.5,2.0) {$y = \ln x$};
\end{axis}
\end{tikzpicture}
\caption{The natural logarithm. It is defined only for $x > 0$, it is negative
  for $0 < x < 1$, zero at $x = 1$, and it grows very slowly: $\ln x$ reaches
  $2$ only at $x = e^2 \approx 7.39$. As $x$ approaches $0$, $\ln x$ plunges
  towards $-\infty$.}
\label{fig:ln}
\end{figure}

\begin{pitfall}[What logarithms do not do]
$\log(x + y)$ is \emph{not} $\log x + \log y$; logarithms turn
\emph{products} into sums, not sums into sums. For example,
$\ln(1 + 1) = \ln 2 \approx 0.693$, but $\ln 1 + \ln 1 = 0$. Also, $\log 0$ and
the logarithm of a negative number are undefined: in code, \texttt{log(0)}
returns $-\infty$ and \texttt{log(-1)} returns ``not a number''.
\end{pitfall}

\begin{example}[Solving exponential equations]
Solve $2^x = 20$ and $e^{2x} = 7$.

\sol Take logarithms of both sides. From $2^x = 20$: $\ln(2^x) = \ln 20$, so
$x\ln 2 = \ln 20$ and $x = \frac{\ln 20}{\ln 2} \approx \frac{2.996}{0.693}
\approx 4.32$. (Check: $2^4 = 16$ and $2^5 = 32$, so the answer must lie between
$4$ and $5$.) From $e^{2x} = 7$: $2x = \ln 7$, so $x = \tfrac12\ln7 \approx
0.973$.
\end{example}

\subsection{Logarithmic scales}

Logarithms measure \emph{orders of magnitude}: $\log_{10}$ of a number is
roughly its number of digits. When values range over many orders of magnitude
we often plot them on a \term{logarithmic scale}, where equal distances
represent equal \emph{ratios}: the distance from $10^{-4}$ to $10^{-3}$ equals
the distance from $10^{-3}$ to $10^{-2}$.

\begin{mlnote}[Why machine learning loves logarithms]
Logarithms are everywhere in machine learning, for three main reasons.
\begin{itemize}
  \item \textbf{Products of probabilities underflow.} The probability that a
    model assigns to a whole data set is a product of many small numbers. With
    $1000$ examples, each given probability $0.01$, the product is
    $10^{-2000}$, far below the smallest number a computer can represent, so
    it becomes $0$. Its logarithm, however, is a perfectly ordinary number:
    $\ln\left(0.01^{1000}\right) = 1000 \ln 0.01 \approx -4605$. So we work
    with sums of log-probabilities (\emph{log-likelihoods}) instead of
    products of probabilities (\cref{ch:statistics}).
  \item \textbf{Loss functions.} The most widely used loss for classification,
    the cross-entropy, is $-\ln p$, where $p$ is the probability the model gave
    to the correct answer (\cref{ch:dl}). It is $0$ when $p = 1$ and grows
    without limit as $p \to 0$.
  \item \textbf{Hyperparameter search.} Learning rates are usually tried on a
    logarithmic scale, e.g.\ $10^{-5}, 10^{-4}, 10^{-3}, 10^{-2}$, because what
    matters is their order of magnitude.
\end{itemize}
\end{mlnote}

\sectionexercises

\begin{exercise}
Evaluate: $\log_2 32$, $\log_{10} 0.001$, $\log_3 81$, $\ln e^5$, $\log_5 1$,
$\log_2 \tfrac18$.
\end{exercise}
\begin{solution}
$5$, $-3$, $4$, $5$, $0$, $-3$.
\end{solution}

\begin{exercise}
(a) Write $\ln\left(\tfrac{x^2 y}{z}\right)$ in terms of $\ln x$, $\ln y$ and
$\ln z$. (b) Simplify $\ln 8 - \ln 2$ and express it as a multiple of $\ln 2$.
\end{exercise}
\begin{solution}
(a) $2\ln x + \ln y - \ln z$. (b) $\ln\tfrac82 = \ln 4 = \ln 2^2 = 2\ln 2$.
\end{solution}

\begin{exercise}
Solve $3^x = 50$ and $5e^{x} = 40$. Give answers to three decimal places.
\end{exercise}
\begin{solution}
$x = \tfrac{\ln 50}{\ln 3} \approx \tfrac{3.912}{1.099} \approx 3.561$. And
$e^x = 8$, so $x = \ln 8 \approx 2.079$.
\end{solution}

\begin{exercise}
A learning rate starts at $0.1$ and is multiplied by $0.9$ after every epoch,
so $\eta_t = 0.1 \times 0.9^t$. After how many epochs does it first drop below
$0.01$?
\end{exercise}
\begin{solution}
We need $0.9^t < 0.1$. Taking logarithms, $t\ln0.9 < \ln0.1$; since
$\ln 0.9 < 0$, dividing by it flips the inequality:
$t > \tfrac{\ln 0.1}{\ln 0.9} \approx 21.85$. So the first such epoch is
$t = 22$. (Check: $0.1\times0.9^{21} \approx 0.0109$ and
$0.1\times0.9^{22} \approx 0.0098$.)
\end{solution}

\begin{exercise}
A model gives probability $0.05$ to each of $1000$ independent examples. Compute
the logarithm of the product of these probabilities, and use it to find the
power of $10$ that the product is roughly equal to.
\end{exercise}
\begin{solution}
$\ln(0.05^{1000}) = 1000\ln0.05 \approx 1000 \times (-2.996) = -2996$.
Converting to base $10$: $\log_{10}(0.05^{1000}) = \tfrac{-2996}{\ln 10}
\approx -1301$, so the product is about $10^{-1301}$, which underflows to $0$
on a computer. The log-probability $-2996$ is harmless.
\end{solution}

%----------------------------------------------------------------------
\section{Composition of functions}\label{sec:composition}
%----------------------------------------------------------------------

Functions can be chained: the output of one becomes the input of the next.

\begin{definition}[Composition]\label{def:composition}
The \term{composition} of $f$ and $g$, written $f \circ g$ (read ``$f$ after
$g$'' or ``$f$ composed with $g$''), is the function
\[
  (f \circ g)(x) = f(g(x)) .
\]
First apply $g$, then apply $f$ to the result.
\end{definition}

\begin{figure}[htbp]
\centering
\begin{tikzpicture}[>={Stealth[length=2.5mm]},
  box/.style={draw,thick,rounded corners,fill=MLdef!10,minimum width=1.6cm,minimum height=1cm}]
  \node[box] (g) at (0,0) {$g$};
  \node[box] (f) at (4,0) {$f$};
  \draw[->,thick] (-2.4,0) node[left] {$x$} -- (g.west);
  \draw[->,thick] (g.east) -- node[above] {$g(x)$} (f.west);
  \draw[->,thick] (f.east) -- (6.4,0) node[right] {$f(g(x))$};
  \draw[decorate,decoration={brace,amplitude=6pt,mirror},thick]
    (-0.9,-0.75) -- (4.9,-0.75) node[midway,below=8pt] {$f \circ g$};
\end{tikzpicture}
\caption{Composition is an assembly line: $x$ goes through $g$ first, and the
  result goes through $f$.}
\label{fig:composition}
\end{figure}

\begin{example}[Order matters]\label{ex:composition}
Let $g(x) = 2x + 1$ and $f(x) = x^2$. Find $f(g(x))$ and $g(f(x))$, and
evaluate both at $x = 3$.

\sol $f(g(x)) = (2x+1)^2$, while $g(f(x)) = 2x^2 + 1$. At $x = 3$:
$f(g(3)) = f(7) = 49$, but $g(f(3)) = g(9) = 19$. In general
$f \circ g \neq g \circ f$.
\end{example}

It is just as important to be able to \emph{decompose} a complicated function
into simple pieces. For example, $h(x) = e^{-x^2}$ is $f(g(x))$ with inner
function $g(x) = -x^2$ and outer function $f(u) = e^u$. In \cref{ch:calculus}
the \emph{chain rule} will compute slopes of compositions piece by piece, and
that rule is the mathematical heart of how neural networks are trained.

\begin{proposition}[Composing linear functions gives a linear function]
\label{prop:compose-linear}
If $f(x) = ax + b$ and $g(x) = cx + d$, then
$f(g(x)) = a(cx + d) + b = (ac)x + (ad + b)$, which is again linear.
\end{proposition}

\begin{mlnote}[Neural networks are compositions]
A neural network with three layers computes
\[
  \text{output} = f_3\bigl(f_2\bigl(f_1(\text{input})\bigr)\bigr)
  = (f_3 \circ f_2 \circ f_1)(\text{input}),
\]
where each \emph{layer} $f_k$ is a fairly simple function. ``Deep'' learning
simply means composing many layers. Each layer consists of a linear part
followed by a simple non-linear function called an \emph{activation function}
(\cref{sec:activation-gallery}). The non-linear part is essential:
\cref{prop:compose-linear} shows that a composition of linear functions is
just one linear function, so without activation functions a deep network
would be no more powerful than a single straight line. We will see the
many-input version of this fact in \cref{ch:linalg1}.
\end{mlnote}

\sectionexercises

\begin{exercise}
Let $f(x) = x + 3$ and $g(x) = x^2$. Compute $f(g(2))$, $g(f(2))$, and
formulas for $f(g(x))$ and $g(f(x))$.
\end{exercise}
\begin{solution}
$f(g(2)) = f(4) = 7$ and $g(f(2)) = g(5) = 25$. In general
$f(g(x)) = x^2 + 3$ and $g(f(x)) = (x+3)^2$.
\end{solution}

\begin{exercise}
Write each function as a composition of simpler functions:
(a) $h(x) = \sqrt{3x + 1}$; (b) $k(x) = \dfrac{1}{1 + e^{-x}}$.
\end{exercise}
\begin{solution}
(a) $h = f \circ g$ with $g(x) = 3x + 1$ and $f(u) = \sqrt u$.
(b) One possible chain: $x \mapsto -x \mapsto e^{-x} \mapsto 1 + e^{-x}
\mapsto \frac{1}{1 + e^{-x}}$, i.e.\ $k = f_4 \circ f_3 \circ f_2 \circ f_1$ with
$f_1(x) = -x$, $f_2(u) = e^u$, $f_3(v) = 1 + v$, $f_4(w) = 1/w$. (This is the
sigmoid function of \cref{sec:activation-gallery}.)
\end{solution}

\begin{exercise}
A ``network'' has one hidden value $h = 2x - 1$ and output $y = -3h + 4$.
(a) Write $y$ as a function of $x$. (b) Now insert a ReLU:
$h = \ReLU(2x - 1)$, $y = -3h + 4$. Compute $y$ for $x = 0, 1, 2$, and decide
whether $y$ is still a linear function of $x$.
\end{exercise}
\begin{solution}
(a) $y = -3(2x - 1) + 4 = -6x + 7$, a linear function.
(b) $x = 0$: $h = \ReLU(-1) = 0$, $y = 4$. $x = 1$: $h = 1$, $y = 1$. $x = 2$:
$h = 3$, $y = -5$. The output changes by $-3$ from $x=0$ to $x=1$ but by $-6$
from $x=1$ to $x=2$; a linear function changes by the same amount for equal
steps, so this is not linear. The ReLU has made the network non-linear.
\end{solution}

%----------------------------------------------------------------------
\section{Inverse functions}\label{sec:inverse-functions}
%----------------------------------------------------------------------

An inverse function undoes what a function does. If $f(x) = 2x + 3$ (``double,
then add $3$''), the inverse is ``subtract $3$, then halve'':
$f^{-1}(y) = \frac{y - 3}{2}$.

\begin{definition}[Inverse function]\label{def:inverse-function}
A function $f\colon A \to B$ has an \term{inverse function}
$f^{-1}\colon B \to A$ if
\[
  f^{-1}(f(x)) = x \text{ for all } x \in A
  \qquad\text{and}\qquad
  f(f^{-1}(y)) = y \text{ for all } y \in B .
\]
\end{definition}

\textbf{How to find an inverse.} Write $y = f(x)$ and solve for $x$ in terms of
$y$. For $y = 2x + 3$: $y - 3 = 2x$, so $x = \frac{y-3}{2}$, which is
$f^{-1}(y)$.

\begin{pitfall}[$f^{-1}$ is not $1/f$]
The $-1$ in $f^{-1}$ is \emph{not} an exponent. For $f(x) = 2x + 3$,
$f^{-1}(y) = \frac{y-3}{2}$, whereas $\frac{1}{f(x)} = \frac{1}{2x+3}$; these
are completely different functions.
\end{pitfall}

\subsection{When does an inverse exist?}

To undo $f$ we must be able to tell, from the output, which input produced
it. So $f$ must never send two different inputs to the same output; such a
function is called \term{one-to-one}. The function $f(x) = x^2$ on $\R$ is not
one-to-one, because $f(2) = f(-2) = 4$: from the output $4$ we cannot tell
whether the input was $2$ or $-2$. But if we \emph{restrict the domain} to
$[0, \infty)$, the function becomes one-to-one and its inverse is the square
root, $f^{-1}(y) = \sqrt{y}$. (Graphically: a function is one-to-one when every
horizontal line meets its graph at most once.)

\subsection{Graphs of inverse functions}

If $(a, b)$ is on the graph of $f$, then $f(a) = b$, so $f^{-1}(b) = a$ and
$(b, a)$ is on the graph of $f^{-1}$. Swapping the coordinates of every point
reflects the graph in the diagonal line $y = x$ (\cref{fig:inverse}).

\begin{figure}[htbp]
\centering
\begin{tikzpicture}
\begin{axis}[mlplot,width=8cm,height=8cm,xmin=-3.2,xmax=4.5,ymin=-3.2,ymax=4.5,
  xlabel=$x$,ylabel=$y$,xtick={-3,...,4},ytick={-3,...,4}]
  \addplot[PlotGray,dashed,domain=-3.2:4.5]{x};
  \addplot[PlotBlue,domain=-3.2:1.46,samples=100]{exp(x)};
  \addplot[PlotRed,domain=0.04:4.5,samples=200]{ln(x)};
  \node[font=\small,PlotBlue,anchor=west] at (axis cs:1.6,4.2) {$y = e^x$};
  \node[font=\small,PlotRed,anchor=north west] at (axis cs:2.6,0.85) {$y = \ln x$};
  \node[font=\small,PlotGray,anchor=north east] at (axis cs:4.4,3.4) {$y = x$};
  \addplot[only marks,mark=*,mark size=2pt,PlotBlue] coordinates {(0,1) (1,2.71828)};
  \addplot[only marks,mark=*,mark size=2pt,PlotRed] coordinates {(1,0) (2.71828,1)};
\end{axis}
\end{tikzpicture}
\caption{$e^x$ and $\ln x$ are inverse functions, so their graphs are mirror
  images in the line $y = x$. The point $(0, 1)$ on $e^x$ corresponds to
  $(1, 0)$ on $\ln x$, and $(1, e)$ to $(e, 1)$.}
\label{fig:inverse}
\end{figure}

\begin{theorem}[Exponential and logarithm are inverses]
\label{thm:exp-ln-inverse}
For every real $x$, $\ln(e^x) = x$, and for every $x > 0$, $e^{\ln x} = x$.
The same holds for any base: $\log_a(a^x) = x$ and $a^{\log_a x} = x$.
\end{theorem}

This is just the definition of the logarithm restated, but it is extremely
useful for solving equations, as we saw in \cref{sec:logarithms}.

\begin{mlnote}[Sigmoid and logit]
The \term{sigmoid function}
\[
  \sigma(z) = \frac{1}{1 + e^{-z}}
\]
(the Greek letter $\sigma$ is ``sigma'') squashes any real number $z$ into the
interval $(0, 1)$, so its output can be read as a probability:
$\sigma(0) = 0.5$, $\sigma(2) \approx 0.881$, $\sigma(-2) \approx 0.119$. Its
inverse is the \term{logit} function,
\[
  \operatorname{logit}(p) = \ln\frac{p}{1 - p}, \qquad 0 < p < 1 .
\]
To check, solve $p = \frac{1}{1+e^{-z}}$ for $z$: $1 + e^{-z} = \frac1p$, so
$e^{-z} = \frac{1-p}{p}$, so $-z = \ln\frac{1-p}{p}$, i.e.\
$z = \ln\frac{p}{1-p}$. The quantity $\frac{p}{1-p}$ is the \emph{odds} (a
probability of $0.8$ means odds of $4$ to $1$), so the logit is the
``log-odds''. In logistic regression and neural network classifiers, the raw
scores $z$ before the sigmoid are called \emph{logits} for this reason.
\end{mlnote}

\sectionexercises

\begin{exercise}
Find the inverse of each function and check your answer:
(a) $f(x) = 5x - 2$; (b) $g(x) = \frac{x - 3}{4}$; (c) $h(x) = x^3 + 1$.
\end{exercise}
\begin{solution}
(a) $y = 5x - 2 \Rightarrow x = \tfrac{y + 2}{5}$, so
$f^{-1}(y) = \tfrac{y+2}{5}$. Check: $f^{-1}(f(x)) = \tfrac{5x - 2 + 2}{5} = x$.
(b) $g^{-1}(y) = 4y + 3$. (c) $y = x^3 + 1 \Rightarrow x = \sqrt[3]{y - 1}$, so
$h^{-1}(y) = \sqrt[3]{y-1}$.
\end{solution}

\begin{exercise}
Is $f(x) = x^2 + 1$ one-to-one on $\R$? On $[0, \infty)$? Find the inverse in
the second case, including its domain.
\end{exercise}
\begin{solution}
Not on $\R$, since $f(1) = f(-1) = 2$. On $[0, \infty)$ it is one-to-one: from
$y = x^2 + 1$ with $x \ge 0$ we get $x = \sqrt{y - 1}$. So
$f^{-1}(y) = \sqrt{y-1}$, defined for $y \ge 1$ (the range of $f$).
\end{solution}

\begin{exercise}
Compute $\sigma(0)$, $\operatorname{logit}(0.5)$ and
$\operatorname{logit}(0.9)$. Then check that the sigmoid undoes the logit, by
computing $\sigma(\operatorname{logit}(0.9))$.
\end{exercise}
\begin{solution}
$\sigma(0) = \tfrac{1}{1+1} = 0.5$; $\operatorname{logit}(0.5) = \ln 1 = 0$;
$\operatorname{logit}(0.9) = \ln\tfrac{0.9}{0.1} = \ln 9 \approx 2.197$. Then
$\sigma(\ln 9) = \tfrac{1}{1 + e^{-\ln 9}} = \tfrac{1}{1 + 1/9} =
\tfrac{9}{10} = 0.9$.
\end{solution}

\begin{exercise}
Data are often standardised with $z = \frac{x - \mu}{s}$, where $\mu = 50$ and
$s = 10$. Find the inverse transformation. A model trained on standardised
data predicts $z = 1.3$; what is the prediction on the original scale?
\end{exercise}
\begin{solution}
Solving for $x$: $x = \mu + sz = 50 + 10z$. For $z = 1.3$, $x = 63$.
\end{solution}

%----------------------------------------------------------------------
\section{How parameters move graphs}\label{sec:transformations}
%----------------------------------------------------------------------

Machine learning models contain parameters, and changing a parameter changes
the graph of the model. It is very useful to be able to predict how. Starting
from any function $f$ and numbers $c > 0$:

\begin{center}
\renewcommand{\arraystretch}{1.25}
\begin{tabular}{@{}ll@{}}
\toprule
\textbf{New function} & \textbf{Effect on the graph of $f$}\\
\midrule
$f(x) + c$ & shift \emph{up} by $c$\\
$f(x) - c$ & shift \emph{down} by $c$\\
$f(x - c)$ & shift \emph{right} by $c$\\
$f(x + c)$ & shift \emph{left} by $c$\\
$c\,f(x)$ & stretch vertically by the factor $c$\\
$f(cx)$ & squeeze horizontally by the factor $c$ (if $c > 1$)\\
$-f(x)$ & reflect in the $x$-axis (upside down)\\
$f(-x)$ & reflect in the $y$-axis (left--right)\\
\bottomrule
\end{tabular}
\end{center}

\begin{pitfall}[Horizontal shifts go the ``wrong'' way]
$f(x - 3)$ moves the graph $3$ units to the \emph{right}, not the left. Reason:
the new function takes at $x = 3$ the value that $f$ took at $0$, so every
feature of the graph now happens $3$ units later. We saw this with
$(x - 3)^2 + 1$, whose vertex is at $x = +3$ (\cref{fig:parabolas}).
\end{pitfall}

\begin{mlnote}[What the weight and bias of a neuron do]
A single artificial neuron with one input computes $\sigma(wx + b)$. Writing
$wx + b = w\left(x + \tfrac bw\right)$ shows that:
\begin{itemize}
  \item the weight $w$ squeezes the sigmoid horizontally: larger $\abs{w}$ gives
    a steeper, more switch-like transition, and a negative $w$ flips it;
  \item the bias $b$ shifts the graph horizontally: the transition (where the
    output is $0.5$) happens where $wx + b = 0$, i.e.\ at $x = -\tfrac bw$.
\end{itemize}
\Cref{fig:sigmoid-family} shows this. By choosing $w$ and $b$, the neuron can
place a soft ``switch'' wherever it is needed; combining many such switches
lets a network build very complicated functions.
\end{mlnote}

\begin{figure}[htbp]
\centering
\begin{tikzpicture}
\begin{axis}[mlplot,width=11cm,height=6.5cm,xmin=-6,xmax=6,ymin=-0.1,ymax=1.15,
  xlabel=$x$,ylabel=$y$,xtick={-6,-4,...,6},ytick={0,0.5,1},
  legend style={font=\small,at={(0.5,-0.1)},anchor=north,legend columns=2,
    draw=none,/tikz/every even column/.append style={column sep=5mm}},
  legend cell align=left]
  \addplot[PlotBlue,domain=-6:6,samples=150]{1/(1+exp(-x))};
  \addlegendentry{$\sigma(x)$: $w=1$, $b=0$}
  \addplot[PlotRed,domain=-6:6,samples=150]{1/(1+exp(-4*x))};
  \addlegendentry{$\sigma(4x)$: steeper}
  \addplot[PlotGreen,domain=-6:6,samples=150]{1/(1+exp(-(x-3)))};
  \addlegendentry{$\sigma(x - 3)$: shifted right}
  \addplot[PlotOrange,domain=-6:6,samples=150]{1/(1+exp(x))};
  \addlegendentry{$\sigma(-x)$: flipped}
  \addplot[PlotGray,dashed,domain=-6:6]{0.5};
\end{axis}
\end{tikzpicture}
\caption{The sigmoid $\sigma(wx + b)$ for different weights and biases. The
  weight controls steepness and direction; the bias controls where the
  transition (output $0.5$, dashed line) happens.}
\label{fig:sigmoid-family}
\end{figure}

\sectionexercises

\begin{exercise}
Describe how to obtain the graph of $g(x) = (x + 1)^2 - 4$ from the graph of
$x^2$. Where is its vertex? Where are its roots?
\end{exercise}
\begin{solution}
Shift $x^2$ one unit to the left and four units down. The vertex is at
$(-1, -4)$. Roots: $(x+1)^2 = 4$, so $x + 1 = \pm2$, i.e.\ $x = 1$ or $x = -3$.
\end{solution}

\begin{exercise}
Starting from $f(x) = \abs{x}$, write down the function whose graph is shifted
$3$ units to the right and stretched vertically by a factor of $2$.
\end{exercise}
\begin{solution}
$g(x) = 2\abs{x - 3}$. Its graph is a V with its tip at $(3, 0)$ and slopes
$\pm2$.
\end{solution}

\begin{exercise}
A neuron computes $\sigma(4x - 2)$. At which input is its output $0.5$? Is it
steeper or less steep than $\sigma(x)$? What changes if the weight becomes
$-4$ (with the same bias)?
\end{exercise}
\begin{solution}
The output is $0.5$ when $4x - 2 = 0$, i.e.\ $x = 0.5$. It is four times as
steep as $\sigma(x)$. With weight $-4$, the function $\sigma(-4x - 2)$
decreases instead of increasing, and its midpoint moves to $x = -0.5$.
\end{solution}

\begin{exercise}
Describe the graphs of $e^{x-2}$, $e^{-x}$, $-e^x$ and $e^x - 1$ in terms of
the graph of $e^x$.
\end{exercise}
\begin{solution}
$e^{x-2}$: shifted $2$ to the right. $e^{-x}$: reflected in the $y$-axis (a
decaying exponential). $-e^x$: reflected in the $x$-axis (always negative).
$e^x - 1$: shifted down by $1$, so it passes through the origin and
approaches $-1$ for very negative $x$.
\end{solution}

%----------------------------------------------------------------------
\section{A gallery of functions used in machine learning}
\label{sec:activation-gallery}
%----------------------------------------------------------------------

\subsection{Piecewise functions}

Some functions are defined by different formulas on different parts of the
domain; these are called \term[piecewise function]{piecewise functions}. We
have already met two: the absolute value $\abs{x}$ (\cref{def:absolute-value})
and $\ReLU(x) = \max(0, x)$, which equals $0$ for $x < 0$ and $x$ for
$x \ge 0$. Another is the \term{step function},
\[
  \operatorname{step}(x) = \begin{cases} 1 & \text{if } x \ge 0,\\
  0 & \text{if } x < 0. \end{cases}
\]

\subsection{Activation functions}

Neural networks use a small number of non-linear functions, called
\term[activation function]{activation functions}, after each linear step.
\Cref{fig:activations} shows the most important ones.

\begin{figure}[htbp]
\centering
\begin{tikzpicture}
\begin{groupplot}[group style={group size=3 by 2,horizontal sep=1.1cm,vertical sep=1.4cm},
  mlplot,width=5cm,height=4.2cm,xmin=-4,xmax=4,xtick={-4,-2,0,2,4},
  title style={font=\small,yshift=-4pt}]
  \nextgroupplot[title={step},ymin=-0.4,ymax=1.4,ytick={0,1}]
    \addplot[PlotBlue,domain=-4:0] {0};
    \addplot[PlotBlue,domain=0:4] {1};
    \addplot[only marks,mark=*,mark size=1.6pt,PlotBlue] coordinates {(0,1)};
    \addplot[only marks,mark=*,mark size=1.6pt,PlotBlue,mark options={fill=white}] coordinates {(0,0)};
  \nextgroupplot[title={sigmoid $\sigma(x)$},ymin=-0.4,ymax=1.4,ytick={0,0.5,1}]
    \addplot[PlotBlue,domain=-4:4,samples=100] {1/(1+exp(-x))};
  \nextgroupplot[title={$\tanh(x)$},ymin=-1.4,ymax=1.4,ytick={-1,0,1}]
    \addplot[PlotBlue,domain=-4:4,samples=100] {tanh(x)};
  \nextgroupplot[title={$\ReLU(x)$},ymin=-1.5,ymax=4.2,ytick={0,2,4}]
    \addplot[PlotBlue,domain=-4:4,samples=100] {max(0,x)};
  \nextgroupplot[title={leaky ReLU ($\alpha = 0.1$)},ymin=-1.5,ymax=4.2,ytick={0,2,4}]
    \addplot[PlotBlue,domain=-4:4,samples=100] {max(0.1*x,x)};
  \nextgroupplot[title={softplus $\ln(1+e^x)$},ymin=-1.5,ymax=4.2,ytick={0,2,4}]
    \addplot[PlotBlue,domain=-4:4,samples=100] {ln(1+exp(x))};
    \addplot[PlotGray,dashed,domain=-4:4,samples=100] {max(0,x)};
\end{groupplot}
\end{tikzpicture}
\caption{Common activation functions. The dashed line in the last panel is
  the ReLU, which softplus smooths out.}
\label{fig:activations}
\end{figure}

\begin{itemize}
  \item \textbf{Step.} Used in the earliest artificial neuron, the perceptron
    (1958): the neuron either fires ($1$) or not ($0$). Its graph has a jump,
    which (as we will see in \cref{ch:calculus}) makes it useless for
    gradient-based training.
  \item \textbf{Sigmoid}, $\sigma(x) = \frac{1}{1 + e^{-x}}$. A smooth
    version of the step, with outputs in $(0, 1)$. Still used to output
    probabilities for yes/no questions.
  \item \textbf{Hyperbolic tangent},
    $\tanh(x) = \frac{e^x - e^{-x}}{e^x + e^{-x}}$. Shaped like the sigmoid but
    with outputs in $(-1, 1)$ and centred at $0$. In fact
    $\tanh(x) = 2\sigma(2x) - 1$.
  \item \textbf{ReLU}, $\max(0, x)$. Simple and fast, and the default choice in
    modern networks since about 2012.
  \item \textbf{Leaky ReLU}, $\max(\alpha x, x)$ for a small $\alpha$ such as
    $0.01$ or $0.1$: like ReLU but with a small slope for negative inputs, so
    that the neuron never becomes completely ``dead''.
  \item \textbf{Softplus}, $\ln(1 + e^x)$: a smooth version of ReLU. Modern
    networks also use other smooth variants of ReLU with names like GELU and
    SiLU, which look similar.
\end{itemize}

\begin{mlnote}[Why so many activation functions?]
All activation functions do the same basic job: they bend the otherwise
straight-line behaviour of a layer, so that compositions of layers can
represent curved, complicated functions. They differ in their ranges
(probabilities need $(0,1)$), in whether they are smooth, and above all in how
easily the network can be trained through them, a question about their
slopes that we will understand in \cref{ch:calculus,ch:dl}.
\end{mlnote}

\sectionexercises

\begin{exercise}
Compute $\sigma(0)$, $\tanh(0)$, $\ReLU(0)$ and $\operatorname{softplus}(0)$.
\end{exercise}
\begin{solution}
$\sigma(0) = \tfrac12$; $\tanh(0) = \tfrac{1 - 1}{1 + 1} = 0$; $\ReLU(0) = 0$;
$\operatorname{softplus}(0) = \ln(1 + 1) = \ln 2 \approx 0.693$.
\end{solution}

\begin{exercise}
Show that $\sigma(-x) = 1 - \sigma(x)$ for every $x$. Interpret this for a
classifier whose output $\sigma(z)$ is the probability of the class ``yes''.
\end{exercise}
\begin{solution}
$1 - \sigma(x) = 1 - \frac{1}{1 + e^{-x}} = \frac{e^{-x}}{1 + e^{-x}}$.
Multiplying top and bottom by $e^x$ gives $\frac{1}{e^{x} + 1} = \sigma(-x)$.
Interpretation: if the logit $z$ gives probability $\sigma(z)$ for ``yes'', then
the probability of ``no'', $1 - \sigma(z)$, is obtained by flipping the sign of
the logit.
\end{solution}

\begin{exercise}
Show that $\tanh(x) = 2\sigma(2x) - 1$.
\end{exercise}
\begin{solution}
$2\sigma(2x) - 1 = \frac{2}{1 + e^{-2x}} - 1 = \frac{1 - e^{-2x}}{1 + e^{-2x}}$.
Multiplying numerator and denominator by $e^{x}$ gives
$\frac{e^x - e^{-x}}{e^x + e^{-x}} = \tanh(x)$.
\end{solution}

\begin{exercise}
Compute the leaky ReLU with $\alpha = 0.01$ at $x = -5$, $0$ and $3$. Then
compute $\operatorname{softplus}(5)$ and $\operatorname{softplus}(-5)$ and compare
with $\ReLU(5)$ and $\ReLU(-5)$.
\end{exercise}
\begin{solution}
Leaky ReLU: $\max(-0.05, -5) = -0.05$, $0$, and $3$.
$\operatorname{softplus}(5) = \ln(1 + e^5) \approx 5.0067$ and
$\operatorname{softplus}(-5) = \ln(1 + e^{-5}) \approx 0.0067$, very close to
$\ReLU(5) = 5$ and $\ReLU(-5) = 0$. Softplus differs noticeably from ReLU only
near $0$.
\end{solution}

%----------------------------------------------------------------------
\section{Chapter summary}\label{sec:ch2-summary}
%----------------------------------------------------------------------

\begin{keytakeaways}
\begin{itemize}
  \item A function assigns to each input exactly one output; it has a domain,
    a codomain and a range. A machine learning model is a function, and
    learning is a search for a good function.
  \item The graph of $f$ is the set of points $(x, f(x))$; roots, intercepts,
    increasing/decreasing behaviour and extreme values can be read from it.
  \item Linear functions $mx + b$ have constant slope $m$ (rise over run, a rate
    of change) and intercept $b$. In ML they are written $wx + b$.
  \item Quadratics $ax^2 + bx + c$ have parabolic graphs with vertex at
    $x = -\frac{b}{2a}$; roots come from factoring or the quadratic formula.
    Squared-error losses are quadratics, minimised by the mean.
  \item Exponentials $a^x$ model multiplicative growth and decay; the natural
    base is $e \approx 2.718$. Logarithms undo exponentials and turn products
    into sums, which is why ML works with log-probabilities.
  \item Composition $f(g(x))$ chains functions; neural networks are long
    compositions. Linear functions composed stay linear, so non-linear
    activations are essential.
  \item Inverse functions undo functions; $e^x$ and $\ln x$ are inverses, and
    so are the sigmoid and the logit.
  \item $f(x - c)$ shifts right, $f(x) + c$ shifts up, $f(cx)$ squeezes. In a
    neuron $\sigma(wx + b)$, $w$ sets the steepness and $b$ the location.
  \item Key activation functions: step, sigmoid, tanh, ReLU, leaky ReLU,
    softplus.
\end{itemize}
\end{keytakeaways}

So far all our functions had one input. Real data has many features, and a
data point is best thought of as a point in a space of many dimensions. The
next chapter builds the geometric intuition we need for this: coordinates,
distances, angles and trigonometry.
